
\section*{初版にむけたあとがき}

このテキストは専修大学人間科学部心理学専攻2年次配当の講義，心理学データ解析応用のテキストです。1年次に姉妹講義であるデータ解析基礎を履修済みであることを想定していますが，単位未習得であってもわかるように心がけて書いているつもりです。

2020年から始まったコロナの問題で，2020年度の前期はオンライン，後期は一部対面での授業となりました。翌2021年度は逆に前期は対面，後期はオンラインが主の講義スタイルになりました。それまでは口頭でできていたことが，より明示的に静的な資料として事前に提示されている必要があると考え，授業内容の書き起こしをしていくことになりました。
データ解析基礎と同様，詳細なシラバスを準備してそれに対応するような授業教材を作ることにしました。基礎編と併せてご利用いただければと思います。
なお，付録など一部基礎編と内容が重複するところがあります。

前半は多変量解析を，後半はベイジアンモデリングという2つのテーマがあるため，どちらかだけ履修することもあるでしょう。とはいえ，行列表現やRの使い方など，両者を通じての説明もありますからテキストとしては一冊にまとめました。



このテキストにはバージョン番号が付与されています。今回はバージョン2.0.1となっています。バージョン番号は，メジャー，マイナー，パッチの3つの要素から構成されています。
今はメジャーバージョン2です。
シラバスの骨組みが変われば，この数字が1つ増えることになります。
チャプターレベルの変更があればマイナーの数字が上がります。
誤字脱字など，細かな変更はパッチの数字が変わります。
なにせ自作のテキストですから，誤字脱字や誤変換，数字や参照のミスが多く，受講生に指摘してもらってはアップデートする，ということを繰り返してここまでやってきました。


最新版のシラバスやテキストは，サイト\footnote{\url{https://kosugitti.github.io/psychometrtics_syllabus/}}で公開されています。
GitHubを使って管理し，図版のほとんどはRを使って作っています。版組にはもちろんLua\LaTeX を使っています。
バージョンがこまめに上がるものですし，リンク機能などPDFならではの利点もあるので，できれば電子デバイスで電子ファイルのまま利用してもらえればと思っています。
しかし受講生の多くはプリントアウトして使っているようですから，この度実費で製本したものが入手できるよう，Kindle Direct Publishingを利用してペーパーバック版も用意しました。
専修大学人間科学部心理学専攻の学生向けに用意したテキストではありますが，どなたでもお読みいただけるようになっています。内容についてはできるだけの注意を払って，誤解や間違いがないよう気をつけてはおりますが，おかしなところがありましたら著者までご連絡ください。

また原稿そのものはCreative Commons BY-SA(CC BY-SA)ライセンスVersion 4.0に基づいて提供しています。加筆修正や再頒布など，ライセンスに従ってご活用いただければと思います。

\begin{flushright}
    \date{\today}\\
    小杉考司
\end{flushright}

